\chapter{Conclusion}

In this thesis,
we have investigated the design of practical threshold key generation and signing schemes.
We have presented several constructions which minimize barriers to implementation,
including reducing network latency,
communication overhead,
and limiting state-keeping.

In Chapter~\ref{chp:frost},
we present FROST,
a Flexible Round-Optimized Schnorr Threshold Signature.
We show how FROST can be used either as a two-round protocol,
or optimized to a single-round online signing protocol,
with a batched preprocessing phase.
In addition,
we present a two-round distributed key generation scheme,
that builds upon Pedersen verifiable secret sharing,
and Schnorr proofs of knowledge.

In Chapter~\ref{chp:storm},
we present new game-based definitions for distributed key generation schemes.
We then present a generic construction that employs as building blocks a non-interactive key-exchange (NIKE) scheme,
an aggregated verifiable secret sharing (AgVSS) scheme,
and a secure hash function.
Finally,
we present Storm,
a concrete instantiation of our generic construction that is secure assuming the hardness of the Computational Diffie-Hellman (CDH) assumption.

Finally, in Chapter~\ref{chp:arctic},
we present Arctic,
a deterministic and stateless threshold Schnorr signature scheme.
Arctic builds upon pseudorandom secret sharing as a building block,
and allows all participants to perform a two-round threshold signing operation,
without requiring participants to maintain state in between or across signing
rounds.
Arctic requires the assumption of an honest minority,
and is efficient for small signing groups (i.e., when the total number of
possible signers is less than $25$).

\section{Future Work}

While this thesis presents advancements in the design of practical threshold key generation and signatures,
many problems remain.
For example, this thesis largely considers the problem of efficient and practical threshold Schnorr signatures.
However, it is possible that the techniques employed in this work are equally applications in the setting of other signature schemes,
such as threshold ECDSA.
While randomized threshold ECDSA signatures under honest-majority assumptions exist that similarly build upon pseudorandom secret sharing~\cite{DamgardJNPO22,DoernerKLS23,DalskovOKSS20},
the techniques we employ in FROST may likewise potentially be applicable to improving the efficiency of threshold ECDSA signatures.

Another avenue for future research is the investigation of the minimal security
assumptions required for two-round threshold Schnorr signature schemes.
While Arctic can be proven under standard assumptions,
FROST requires the stronger Algebraic One-More Discrete Logarithm (AOMDL)
assumption.
It is interesting to consider whether this difference is inherent,
or if it is possible to prove two-round randomized Schnorr threshold signature
secure under weaker assumptions.
